\begin{ouverture}[courbe]{9}{Fonctions linéaires et affines}{Partie 2 · Proportionnalité et fonctions}
{Proportionnalité · Droites · Coefficients · Inéquations · Tarifs}
\objectif{Reconnaître une fonction linéaire ou affine.}
\objectif{Tracer une droite et déterminer graphiquement ses coefficients.}
\objectif{Comparer des modèles et résoudre des inéquations par le calcul et le graphique.}
\prerequis{Fonctions, équations et inéquations du premier degré.}
\end{ouverture}
\pagestyle{otMain}\essentiel\fontsize{11}{15}\selectfont
\essSection{1}{Deux familles de fonctions}
\begin{definition}Une fonction linéaire est une fonction de la forme $f:x\mapsto ax$. Elle modélise une situation de proportionnalité, de coefficient $a$, et sa représentation est une droite passant par l’origine.\end{definition}
\begin{definition}Une fonction affine est une fonction de la forme $g:x\mapsto ax+b$. Sa représentation est une droite ; $b=g(0)$ est son ordonnée à l’origine. Le coefficient $a$, appelé coefficient directeur ou pente, indique de combien varie l’ordonnée lorsque l’abscisse augmente de 1.\end{definition}
\begin{exemple}Pour $g(x)=2x-3$, le point $(0\,;-3)$ est sur la droite. Quand l’abscisse augmente de 1, l’ordonnée augmente de 2. Les points $(0\,;-3)$ et $(2\,;1)$ suffisent pour tracer la droite.\end{exemple}
\begin{propriete}Un coefficient $a$ positif correspond à une droite montante de gauche à droite, un coefficient négatif à une droite descendante. Si $a=0$, la fonction est constante.\end{propriete}
\begin{attention}« Affine » n’implique pas « proportionnel ». Il faut en plus $b=0$ pour obtenir une fonction linéaire.\end{attention}
\clearpage
\essSection{2}{Retrouver une formule et comparer}
\begin{savoirfaire}[Coefficients sur un graphique]Lire $b$ à l’intersection avec l’axe des ordonnées. Pour $a$, partir d’un point de la droite, avancer d’une unité vers la droite et lire la variation verticale, positive en montant, négative en descendant.\end{savoirfaire}
\begin{savoirfaire}[À partir de deux points]Pour deux points de la droite $(x_1\,;y_1)$ et $(x_2\,;y_2)$ avec $x_1\ne x_2$ :
\[a=\frac{y_2-y_1}{x_2-x_1}.\]
Utiliser ensuite $y_1=ax_1+b$ pour retrouver $b$. Une fonction linéaire est déterminée par un seul nombre non nul et son image.\end{savoirfaire}
\begin{savoirfaire}[Inéquations et seuils]Pour comparer deux tarifs, résoudre leur égalité pour repérer le seuil, puis l’inéquation correspondante, par le calcul ou en comparant les positions des droites. Revenir au nombre entier d’unités lorsque la situation l’exige.\end{savoirfaire}
\begin{attention}Une droite verticale ne représente pas une fonction : la même entrée aurait plusieurs images. Une formule affine doit être justifiée par le modèle ou les données.\end{attention}
\clearpage\rubrique{Découverte}
\begin{blocLivre}{ex-9-1}\begin{exercice}[1]{Reconnaître les formes}
Parmi $f(x)=-3x$, $g(x)=4x+1$, $h(x)=x^2$ et $k(x)=7$, préciser les fonctions affines et les fonctions linéaires.
\end{exercice}\end{blocLivre}
\begin{blocLivre}{ex-9-2}\begin{exercice}[1]{Calculer avec une linéaire}
Pour $f(x)=-2{,}5x$, calculer $f(4)$ et l’antécédent de 15.
\end{exercice}\end{blocLivre}
\begin{blocLivre}{ex-9-3}\begin{exercice}[1]{Lire les coefficients}
Pour $g(x)=-3x+8$, donner le coefficient directeur, l’ordonnée à l’origine et $g(2)$.
\end{exercice}\end{blocLivre}
\begin{blocLivre}{ex-9-4}\begin{exercice}[1]{Tracer une linéaire}
Tracer la représentation de $f(x)=1{,}5x$ à l’aide de deux points.
\label{fig-ex-9-4}\figureGeo{affine-trace-lineaire}
\end{exercice}\end{blocLivre}
\begin{blocLivre}{ex-9-5}\begin{exercice}[1]{Tracer une affine}
Tracer $g(x)=-x+4$ à l’aide de deux points, puis lire son antécédent de 0.
\label{fig-ex-9-5}\figureGeo{affine-trace}
\end{exercice}\end{blocLivre}

\clearpage\rubrique{Entraînement}
\begin{blocLivre}{ex-9-6}\begin{exercice}[2]{Une linéaire à retrouver}
Une fonction linéaire $f$ vérifie $f(6)=21$. Donner son expression et $f(-2)$.
\end{exercice}\end{blocLivre}
\begin{blocLivre}{ex-9-7}\begin{exercice}[2]{Deux points, une formule}
Une fonction affine passe par $A(2\,;1)$ et $B(5\,;10)$. Déterminer son expression.
\label{fig-ex-9-7}\figureGeo{affine-deux-points}
\end{exercice}\end{blocLivre}
\begin{blocLivre}{ex-9-8}\begin{exercice}[2]{Retrouver un tarif}
Un service affine facture 14 € pour 2 heures et 26 € pour 5 heures. Déterminer le forfait fixe et le prix par heure.
\end{exercice}\end{blocLivre}
\begin{blocLivre}{ex-9-9}\begin{exercice}[2]{Coefficients lus sur un graphique}
Déterminer graphiquement les coefficients $a$ et $b$ de la fonction affine représentée, puis donner son expression.
\label{fig-ex-9-9}\figureGeo{affine-lecture}
\end{exercice}\end{blocLivre}
\begin{blocLivre}{ex-9-10}\begin{exercice}[2]{Comparer par une différence}
On donne $f(x)=4x-1$ et $g(x)=4x+6$. Comparer leurs images pour toute entrée.
\label{fig-ex-9-10}\figureGeo{affine-paralleles}
\end{exercice}\end{blocLivre}

\clearpage\rubrique{Problèmes}
\begin{blocLivre}{ex-9-11}\begin{exercice}[3]{Deux abonnements}
Pour $n$ séances, l’offre A coûte $20+3n$ euros et l’offre B $8+5n$ euros. Déterminer le nombre de séances donnant le même coût, puis résoudre l’inéquation $20+3n<8+5n$ et conclure.
\label{fig-ex-9-11}\figureGeo{affine-tarifs}
\end{exercice}\end{blocLivre}
\begin{blocLivre}{ex-9-12}\begin{exercice}[3]{Une vidange à débit constant}
Une cuve contient 900 L. Elle perd 35 L par minute. Donner le volume restant $V(t)$ avant qu’elle soit vide. Quand reste-t-il 200 L ? Après 30 minutes, la formule peut-elle encore représenter un volume ?
\end{exercice}\end{blocLivre}
\begin{blocLivre}{ex-9-13}\begin{exercice}[3]{Une vitesse et une inéquation}
Un mobile se déplace à vitesse constante de 5 m/s. On note $d(t)$ la distance parcourue, en mètres, après $t$ secondes. Exprimer $d$, puis déterminer à partir de quel instant la distance dépasse strictement 120 m. Retrouver le résultat graphiquement.
\end{exercice}\end{blocLivre}
\begin{blocLivre}{ex-9-14}\begin{exercice}[3]{Un point suffit-il ?}
On sait seulement qu’une fonction affine vérifie $f(2)=7$. Peut-on déterminer sa formule ? Donner deux formules différentes possibles.
\end{exercice}\end{blocLivre}
\begin{blocLivre}{ex-9-15}\begin{exercice}[3]{Un seuil avec contrainte}
Deux locations pour $d$ jours entiers coûtent $A(d)=45+12d$ et $B(d)=15+18d$. Un budget de 120 € est prévu. Déterminer la meilleure offre pour 6 jours et le nombre maximal de jours possible dans chaque offre.
\end{exercice}\end{blocLivre}

\clearpage
